\section{The actual Segre pullback retains whole original sections and model norms}
\label{sec:projective-segre-pullback-sections}

\begin{definition}[The original pair morphism and its Segre composite]
\label{def:original-pair-segre-morphism}
\lean{EnvelopeProperty.ProjectiveSegrePullbackSections.pair, EnvelopeProperty.ProjectiveSegrePullbackSections.morphism, EnvelopeProperty.ProjectiveSegrePullbackSections.morphism_structure}
\leanok
Let $f:X\to\mathbf P_R^m$ and $g:X\to\mathbf P_R^n$ be actual
original morphisms over the same commutative ring $R$. Their true pair
morphism is the map into the actual fiber product determined by $f,g$.
Composing it with the actual Segre map gives an original morphism
$s:X\to\mathbf P_R^{(m+1)(n+1)-1}$ over the unchanged original base.
\end{definition}
\begin{proof}
The equality of the two original base composites supplies the universal
fiber-product map. Its first and second projections are literally $f$
and $g$. The genuine Segre map has the original first-factor structure
composite, so its composite with this pair map has the original base map
of $f$, which is also that of $g$.
\end{proof}

\begin{theorem}[The derived whole original line isomorphism]
\label{thm:original-segre-pulled-line-iso}
\lean{EnvelopeProperty.ProjectiveSegrePullbackSections.line, EnvelopeProperty.ProjectiveSegrePullbackSections.line_isInvertible, EnvelopeProperty.ProjectiveSegrePullbackSections.lineIso}
\leanok
\uses{def:original-pair-segre-morphism,thm:fc-original-segre-projective-twist}
For these actual original maps there is a genuine whole module
isomorphism
\[
 s^*\mathcal O(1)\simeq f^*\mathcal O(1)\otimes g^*\mathcal O(1).
\]
All three modules are the original invertible pullbacks on $X$.
\end{theorem}
\begin{proof}
The whole original Segre line on the fiber product is the tensor of its
actual factor lines. Pull this proved whole isomorphism back by the
original pair morphism. The actual strong monoidal pullback comparison
identifies its tensor pullback with the tensor of the two pullbacks.
The true factorization comparisons for the pair projections identify
these factors with the original pullbacks along $f$ and $g$. Finally the
true composite pullback comparison identifies its source with the same
original line pulled back along $s$. Composing these genuine whole
isomorphisms gives the assertion. Original projective lines are
invertible and true pullback preserves their invertibility.
\end{proof}

\begin{theorem}[Every actual original coordinate section is retained]
\label{thm:original-segre-pulled-coordinate-section}
\lean{EnvelopeProperty.ProjectiveSegrePullbackSections.coordinateSection, EnvelopeProperty.ProjectiveSegrePullbackSections.lineIso_coordinate}
\leanok
\uses{thm:original-segre-pulled-line-iso,thm:fc-original-segre-projective-twist,lem:fc-original-tensor-section-homs}
Under this constructed whole isomorphism, the pullback of every original
Segre coordinate $Z_{ij}$ is exactly the actual tensor section of the
original pullbacks of $X_i$ and $Y_j$.
\end{theorem}
\begin{proof}
On the actual fiber product, the proved original Segre line isomorphism
takes $Z_{ij}$ to the tensor of the two original pulled coordinate
sections. Apply the actual pair pullback to this equality. Actual
pullback of unit-section morphisms preserves their composition with a
module map. Its genuine monoidal comparison takes the pulled tensor
section to the tensor of the two pulled sections. The actual composite
and factor pullback identities identify these three sections with the
original sections pulled back respectively along $s,f,g$. These are
equalities of whole section morphisms, so no generation or coordinate
matching premise is needed.
\end{proof}

\begin{theorem}[All original whole sections and finite weighted families are isometric]
\label{thm:original-segre-whole-sections-model-isometry}
\lean{EnvelopeProperty.ProjectiveSegrePullbackSections.wholeSectionsEquiv, EnvelopeProperty.ProjectiveSegrePullbackSections.wholeSectionsNorm, EnvelopeProperty.ProjectiveSegrePullbackSections.weightedNormEval}
\leanok
\uses{thm:original-segre-pulled-line-iso,thm:actual-module-iso-global-sections,thm:scheme-global-section-model-norm-isometry,thm:scheme-global-section-weighted-norm-isometry}
The constructed whole line isomorphism induces a genuine equivalence on
the entire original global section spaces. If the original $X$ is proper
over $R$, then at every original valued point over any bounded base
valuation, it preserves the actual proper-model norm of every original
global section. It also preserves the finite weighted logarithmic
evaluation of every original finite family, retaining the same positive
denominator and all real weights, including empty or common-zero families.
\end{theorem}
\begin{proof}
Apply global sections to the constructed whole module isomorphism and
its actual inverse. The two compositions are identities on every whole
section, so this is a genuine equivalence on the full spaces. At every
original proper valued point, the proved whole-module isometry makes
these same section maps norm-preserving. For a finite weighted family,
map its actual sections through this equivalence and retain its weights.
Every term of the finite logarithmic supremum is unchanged by the norm
identity; taking the same finite supremum and multiplying by the same
denominator inverse proves equality of the full weighted evaluations.
This argument includes zero sections and empty families.
\end{proof}

\begin{theorem}[Actual coordinate norms multiply at every original native point]
\label{thm:original-segre-coordinate-model-norm}
\lean{EnvelopeProperty.ProjectiveSegrePullbackSections.coordinateNorm}
\leanok
\uses{thm:original-segre-pulled-coordinate-section,thm:scheme-global-section-tensor-norm,thm:scheme-global-section-model-norm-isometry}
For proper original $X$, the actual proper-model norm of the original
pulled $Z_{ij}$ section equals the product of the actual norms of the
original pulled $X_i$ and $Y_j$ sections, at every original native valued
point. Supports may be nongeneric and either factor norm may be zero.
\end{theorem}
\begin{proof}
The constructed line isomorphism has actual model morphism norm one.
The whole coordinate-section identity therefore identifies the original
Segre coordinate norm with that of the actual tensor section. The true
original tensor norm theorem identifies the latter with the product of
the two original coordinate norms. Both norm theorems retain zero values,
so no nonvanishing or generic-support restriction is required.
\end{proof}
